\section{Proofs for finite betting and subjective probability}
\label{sec:annex-p17-finite-subjective-probability}

\begin{annexitem}{annex:p17.bet-cash-equivalent}{thm:p17-bet-cash-equivalent}{Construction and linearity of the cash equivalent}
\begin{proof}
Additivity gives the forward weak comparison after translation.
If \(x+z\succeq y+z\), apply it with the common vector \(-z\)
to obtain \(x\succeq y\). Apply these equivalences also with
\(x,y\) exchanged. They preserve both indifference and strict preference.
Indifferent objects can be substituted in a weak comparison by
transitivity in each direction.

Suppose \(x\geq_{\rm cw}y\). For every \(\eta>0\), all coordinates
of \(x+\eta\mathbf1\) exceed those of \(y\), so
\(x+\eta\mathbf1\succ y\). If \(x\not\succeq y\), completeness
would give \(y\succ x\). Continuity would then preserve this latter
comparison with \(x\) replaced by \(x+\eta\mathbf1\) for
sufficiently small positive \(\eta\), a contradiction. Thus
\(x\succeq y\). Also, \(t>s\) implies
\(t\mathbf1\succ s\mathbf1\) directly from strict improvement
in every coordinate. Comparisons of sure payments are therefore
exactly comparisons of their real amounts.

For a given \(x\), put \(a=\min_i x_i\) and \(b=\max_i x_i\).
Coordinatewise weak monotonicity gives
\(b\mathbf1\succeq x\succeq a\mathbf1\).
If \(a=b\), the bet is itself that certain vector. Otherwise define
\[
                   H=\{t\in[a,b]:x\succeq t\mathbf1\}.
\]
This set is nonempty and bounded. It is closed relative to \([a,b]\):
if \(t\notin H\), completeness gives \(t\mathbf1\succ x\),
and continuity preserves this comparison for nearby sure amounts.
The map \(t\mapsto t\mathbf1\) is continuous, since
\(d(t\mathbf1,s\mathbf1)=\sqrt K\,|t-s|\).

Let \(c=\sup H\). Closedness gives \(c\in H\), either directly
when \(c=a\), or by taking a sequence of members approaching the
supremum. Thus \(x\succeq c\mathbf1\). If this comparison were
strict, then \(c=b\) would contradict \(b\mathbf1\succeq x\).
For \(c<b\), continuity would give a slightly larger \(t\leq b\)
still satisfying \(x\succ t\mathbf1\), contradicting the supremum.
Hence \(x\sim c\mathbf1\). Two different such amounts are strictly
ordered as sure payments and cannot both be indifferent to \(x\).
This proves existence, uniqueness and the asserted bounds; denote
the amount by \(C(x)\).

The indifferences for \(x,y\) and transitivity give
\[
 x\succeq y\quad\Longleftrightarrow\quad
 C(x)\mathbf1\succeq C(y)\mathbf1
 \quad\Longleftrightarrow\quad C(x)\geq C(y).
\]
For a sure payment, uniqueness gives \(C(t\mathbf1)=t\), in
particular \(C(\mathbf0)=0\). Translate the indifference for
\(x\) by \(y\), and the indifference for \(y\) by
\(C(x)\mathbf1\). Transitivity then gives
\[
 x+y\sim C(x)\mathbf1+y
             \sim(C(x)+C(y))\mathbf1,
\]
so \(C(x+y)=C(x)+C(y)\). Consequently \(C(-x)=-C(x)\),
and repeated addition gives \(C(mx)=mC(x)\) for every integer
\(m\). For a positive integer \(n\),
\(nC((m/n)x)=C(mx)=mC(x)\); thus homogeneity holds for every
rational scalar, including zero and negative values.

To extend it to real scalars, put \(\delta=\|x-y\|_\infty\).
Coordinatewise,
\(y-\delta\mathbf1\leq_{\rm cw}x\leq_{\rm cw}y+\delta\mathbf1\).
Weak monotonicity, representation by \(C\), and its additivity give
\[
                         C(y)-\delta\leq C(x)\leq C(y)+\delta.
\]
This proves the Lipschitz bound. If rationals \(q_n\) converge to
\(\lambda\in\mathbb R\), then
\(\|q_nx-\lambda x\|_\infty=|q_n-\lambda|\|x\|_\infty\to0\).
Therefore \(C(\lambda x)=\lim_n C(q_nx)
=\lim_n q_n C(x)=\lambda C(x)\). No nonzero-coordinate or positive-scalar
restriction is used. All these arguments also apply when \(K=1\).
\end{proof}
\end{annexitem}

\begin{annexitem}{annex:p17.finite-bet-representation}{thm:p17-finite-bet-representation}{The unique normalized linear representation}
\begin{proof}
Under the axioms the preceding lemma gives a linear cash-equivalent
function. Every vector has the finite decomposition
\(x=\sum_{i=1}^K x_i e_i\). Additivity and real homogeneity yield
\[
                          C(x)=\sum_i x_i C(e_i).
\]
Set \(\pi_i=C(e_i)\). Since \(e_i\geq_{\rm cw}\mathbf0\),
weak monotonicity gives \(\pi_i\geq C(\mathbf0)=0\).
Also \(\sum_i e_i=\mathbf1\), so
\(\sum_i\pi_i=C(\mathbf1)=1\). Thus \(\pi\) is a probability
vector, and \(C(x)=\pi\cdot x\) represents the preference.
For uniqueness, any probability vector representing the same preference
must make \(e_i\) indifferent to its stated sure payment \(C(e_i)\).
Its value at that certain vector is the sure amount because its
coordinates sum to one; its value at \(e_i\) is its \(i\)-th
coordinate. Those coordinates must therefore equal \(C(e_i)\)
for every \(i\).

Conversely let \(\pi\) be a probability vector and define the
relation by its linear evaluation. It is complete and transitive by
the real order. Adding a common vector adds the same number
\(\pi\cdot z\) to both evaluations, proving additivity in both
directions. If \(x_i>y_i\) for every \(i\), their finitely many
positive differences have positive minimum \(\eta\). Hence
\[
              \pi\cdot(x-y)\geq\eta\sum_i\pi_i=\eta>0.
\]
This proves the source's weak monotonicity even when some \(\pi_i=0\).

For any two vectors,
\[
            |\pi\cdot(x'-x)|
               \leq\sum_i\pi_i|x_i'-x_i|
               \leq\|x'-x\|_\infty\leq d(x',x).
\]
If \(x\succ y\), let \(g=\pi\cdot x-\pi\cdot y>0\).
Moving each argument by distance less than \(g/4\) leaves the
evaluation gap greater than \(g/2\), proving the required continuity.
When \(K=1\), normalization forces \(\pi_1=1\), giving the
asserted certain-payoff case.
\end{proof}
\end{annexitem}

\begin{annexitem}{annex:p17.bet-halfspaces}{thm:p17-bet-halfspaces}{The separating halfspaces and their orientation}
\begin{proof}
The representation just proved gives
\(x\succeq\mathbf0\) exactly when \(\pi\cdot x\geq0\),
and \(\mathbf0\succ x\) exactly when \(\pi\cdot x<0\).
Thus the displayed sets are disjoint and cover all of \(\mathbb R^K\).
The zero vector belongs to \(\mathcal U\), and
\(-\mathbf1\) belongs to \(\mathcal D\), since their evaluations
are zero and minus one. Both sets are nonempty.

A convex combination of two nonnegative evaluations is nonnegative;
a convex combination of two negative evaluations is negative,
including the endpoint weights. Linearity proves convexity of both
sets. The displayed Lipschitz bound for \(\pi\cdot x\) implies
that a limit of points with nonnegative evaluations still has
nonnegative evaluation, so \(\mathcal U\) is closed. A negative
evaluation persists in a sufficiently small neighborhood, so
\(\mathcal D\) is open.

Normalization gives \(\pi\cdot\mathbf1=1\), so the normal vector
is nonzero, and its coordinates are nonnegative by the representation
proof. Its zero level passes through \(\mathbf0\). At any point
of that level, translations by \(\eta\mathbf1\) and
\(-\eta\mathbf1\), for arbitrarily small \(\eta>0\), have
respectively positive and negative evaluations. Thus the zero level
is the common boundary with the asserted orientation. This supplies
the geometric separation used in the source through the explicit
constructed linear functional.
\end{proof}
\end{annexitem}

\begin{annexitem}{annex:p17.bet-events-null-states}{thm:p17-bet-events-null-states}{Event indicators and the strength of monotonicity}
\begin{proof}
The finite decomposition of the indicator gives
\[
            C(\mathbf1_A)=C\left(\sum_{i\in A}e_i\right)
                 =\sum_{i\in A}C(e_i)=\sum_{i\in A}\pi_i.
\]
This includes the empty event and the whole state set, with respective
values zero and one. These nonnegative normalized masses give exactly
the finite event measure in the definition.

Changing only coordinate \(i\) by a real amount \(t\) changes
the cash equivalent by \(C(x+te_i)-C(x)=t\pi_i\).
If \(\pi_i=0\), every such change is indifferent. Conversely,
if every such change leaves indifference, take \(x=\mathbf0,t=1\)
to obtain \(\pi_i=0\).

If every \(\pi_i>0\) and \(x\geq_{\rm cw}y\) with a strict
inequality in at least one coordinate, all terms in
\(\sum_i\pi_i(x_i-y_i)\) are nonnegative and at least one is
positive. Thus \(x\succ y\). If some \(\pi_i=0\), the comparison
of \(e_i\) and \(\mathbf0\) has the required coordinatewise
improvement but remains an indifference, so the stronger condition
fails. This proves both directions.
\end{proof}
\end{annexitem}

\begin{annexitem}{annex:p17.finite-price-coherence}{thm:p17-finite-price-coherence}{Finite portfolios and the probability axioms}
\begin{proof}
Suppose first that \(q(A)=\sum_{i\in A}\pi_i\) for a probability
vector. For every priced event contract,
\[
               \pi\cdot b_A=\pi\cdot\mathbf1_A
                                  -q(A)\pi\cdot\mathbf1=0.
\]
Finite linearity gives expectation zero for any permitted portfolio,
including one with negative coefficients or no terms. If its payoff
were negative in every state, finiteness and nonemptiness of the state
set would give \(M=\max_i x_i<0\). Its expectation would satisfy
\(\pi\cdot x\leq M\sum_i\pi_i=M<0\), a contradiction.
Thus probability quotes are coherent in the stated sense.

Conversely suppose no strictly negative-everywhere portfolio can be
formed. A nonzero constant portfolio is also impossible: either it
or its negative would be strictly negative everywhere, and negative
quantities are allowed. Since
\(b_\varnothing=-q(\varnothing)\mathbf1\) and
\(b_S=(1-q(S))\mathbf1\), this implies
\(q(\varnothing)=0\) and \(q(S)=1\).
If \(q(A)<0\), then \(b_A=\mathbf1_A-q(A)\mathbf1\) is positive
in every state; its negative is a forbidden portfolio. If \(q(A)>1\),
\(b_A\) itself is negative in every state. Hence
\(0\leq q(A)\leq1\) for every event.

For disjoint events \(A,B\), the pointwise identity
\(\mathbf1_{A\cup B}=\mathbf1_A+\mathbf1_B\) gives
\[
              b_{A\cup B}-b_A-b_B
                   =\bigl(-q(A\cup B)+q(A)+q(B)\bigr)\mathbf1.
\]
This constant portfolio must be zero, so \(q\) is finitely additive
on disjoint events. Iterating over the finitely many singleton states
gives
\(q(A)=\sum_{i\in A}q(\{i\})\).
Set \(\pi_i=q(\{i\})\). They are nonnegative and sum to
\(q(S)=1\), so form the required unique probability vector. On a
finite state set this also gives countable additivity: a disjoint
sequence has at most finitely many nonempty members. Thus the quotes
are a probability measure on all its subsets.

If all coordinates of this vector are positive, a portfolio which
is everywhere nonpositive and negative in at least one state has
strictly negative expectation, again contradicting its zero expectation.
If instead \(\pi_i=0\), then \(q(\{i\})=0\), so
\(-b_{\{i\}}=-e_i\) is an allowed portfolio. It is negative in
that state and zero in all others. A normalized probability vector
with a zero coordinate has at least one other state of positive
mass, so this example does not contradict exclusion of a loss
strict in every state. Both coherence strengths are now distinguished.
\end{proof}
\end{annexitem}

\begin{annexitem}{annex:p17.nonlinear-bet-example}{thm:p17-nonlinear-bet-example}{An additive-background failure under nonlinear utility}
\begin{proof}
The real score \(V\) gives a complete and transitive comparison.
It is a continuous function of the two coordinates, so a positive
score gap persists under sufficiently small changes in both arguments,
as required by the preference continuity definition. The function
\(t\mapsto-e^{-t}\) is strictly increasing. Increasing both
coordinates strictly therefore increases both summands and hence \(V\).

For the specified vectors,
\(V(x)=V(y)=-e^{-1}-1\). But \(x+z=(1,1)\) and
\(y+z=(0,2)\), with
\[
                  V(x+z)-V(y+z)
                    =1+e^{-2}-2e^{-1}=(1-e^{-1})^2>0.
\]
The original indifference includes \(y\succeq x\); additivity would
require \(y+z\succeq x+z\), which the strict reverse comparison
contradicts. Every linear probability-weighted representation satisfies
additivity by the converse already proved, so none represents this
preference on the whole bet space. Nevertheless \(V/2\) is precisely
the equally weighted expected utility of the transformed payoffs,
which proves the distinction stated in the body.
\end{proof}
\end{annexitem}
